\chapter{Particle-Filter Foundations}
\label{ch:bf-particle-filters}

This chapter develops the particle-filter baseline for the differentiable
particle-filter program.  The goal is not yet to make the filter
differentiable.  The goal is to establish, in BayesFilter notation, the
nonlinear filtering recursion, the empirical particle approximation, the
bootstrap particle-filter likelihood estimator, and the exact point at which the
later differentiable program departs from the classical construction.

Particle methods matter because the nonlinear filtering problem is generally not
closed under finite-dimensional moments.  When the predictive law or filtering
law is materially non-Gaussian, a Gaussian closure such as the EKF or UKF may be
a useful approximation but no longer defines the exact filtering object.  The
particle filter instead approximates the filtering law directly as a weighted
random measure.  This makes it the natural baseline for the particle-flow,
proposal-correction, resampling-relaxed, and differentiable transport staircase
that follows.

\section{Objects, assumptions, and claim-status vocabulary}
\label{sec:bf-pf-objects}

The later DPF chapters rely on this chapter as the reference object.  The
following notation is therefore fixed before introducing any differentiable
relaxation.

\begin{center}
\small
\begin{tabular}{p{0.25\textwidth} p{0.65\textwidth}}
\toprule
Object & Meaning in this chapter \\
\midrule
$x_t\in\mathbb{R}^{n_x}$ & latent state at time $t$ \\
$y_t\in\mathbb{R}^{n_y}$ & observation at time $t$ \\
$\theta\in\Theta$ & fixed model parameter while the filtering recursion is
defined \\
$p_\theta(x_t\mid x_{t-1})$ & transition density or Markov kernel density \\
$p_\theta(y_t\mid x_t)$ & observation density, evaluated at the realized
observation $y_t$ \\
$p_\theta(x_t\mid y_{1:t-1})$ & predictive law before assimilating $y_t$ \\
$p_\theta(x_t\mid y_{1:t})$ & filtering law after assimilating $y_t$ \\
$p_\theta(y_t\mid y_{1:t-1})$ & one-step likelihood factor \\
$q_\theta(x_t\mid x_{0:t-1},y_{1:t})$ & sequential proposal density \\
$\tilde w_t^{(i)},w_t^{(i)}$ & unnormalized and normalized particle weights \\
$a_{t-1}^{(i)}$ & ancestor index selected during resampling \\
$\hat\pi_t^N$ & weighted empirical approximation to the filtering law \\
$Z_t^N$ & average incremental bootstrap weight at time $t$ \\
$\hat p_\theta^N(y_{1:T})$ & bootstrap particle-filter estimator of the
marginal likelihood \\
\bottomrule
\end{tabular}
\end{center}

The standing assumptions for the formulas below are the standard ones needed
to make the displayed densities and expectations well defined: the transition
and observation kernels are dominated by reference measures; the proposal has
support wherever the trajectory target has support; the relevant integrals are
finite; and resampling schemes are conditionally unbiased in the sense that the
expected empirical measure after resampling equals the normalized empirical
measure before resampling.  These assumptions are not cosmetic.  They mark the
difference between an exact model identity, a finite-particle Monte Carlo
estimator statement, and an implementation diagnostic.

The chapter uses the following claim-status vocabulary.  The filtering
recursion and likelihood factorization are exact model identities.  Weighted
particles are finite-$N$ Monte Carlo approximations.  The bootstrap likelihood
estimator is an unbiased estimator of the likelihood, not of the log likelihood
and not of the score.  ESS and ancestor-count summaries are diagnostics, not
proofs of posterior correctness or failure.

\section{DPF-block notation and claim-status index}
\label{sec:bf-dpf-block-notation-index}

This chapter begins the DPF block, so the following index fixes the notation
used across the later particle-flow, PF-PF, resampling/OT, learned-OT, HMC, and
debugging chapters.  A symbol may acquire a method-specific subscript in a
later chapter, but its mathematical role should not change silently.

\begin{center}
\scriptsize
\begin{longtable}{@{}p{0.20\textwidth} p{0.24\textwidth} p{0.38\textwidth}@{}}
\caption{DPF-block notation index.}
\label{tab:bf-dpf-notation-index}\\
\toprule
Category & Symbols & Fixed role across the DPF block \\
\midrule
\endfirsthead
\toprule
Category & Symbols & Fixed role across the DPF block \\
\midrule
\endhead
State and observations &
$x_t$, $x_{0:t}$, $y_t$, $y_{1:t}$ &
Latent state, latent trajectory, current observation, and observation history
for the state-space model. \\
Parameters &
$\theta$, $u$ &
Structural/model parameter and, in HMC chapters, unconstrained sampler
coordinate after any parameter transform. \\
Filtering laws &
$p_\theta(x_t\mid y_{1:t-1})$, $p_\theta(x_t\mid y_{1:t})$ &
Predictive and filtering laws; these are the reference statistical objects. \\
Particles and ancestors &
$x_t^{(i)}$, $\tilde x_t^{(j)}$, $a_{t-1}^{(i)}$, $a_j$ &
Pre- or post-update particles, equal-weight output particles, and discrete
ancestor indices. \\
Weights and empirical measures &
$\tilde w_t^{(i)}$, $w_t^{(i)}$, $\hat\pi_t^N$, $\hat\pi_{t,\mathrm{eq}}^N$ &
Unnormalized/normalized weights and weighted/equal-weight empirical measures. \\
Likelihood estimates &
$Z_t^N$, $\hat p_\theta^N(y_{1:T})$, $\widehat L(y\mid\theta,\omega)$ &
Average incremental particle weight, bootstrap likelihood estimator, and
generic randomized likelihood estimator used in pseudo-marginal discussion. \\
Proposals &
$q_\theta$, $q_{t,0}$, $q_{t,1}$, $\gamma_t$ &
Sequential proposal, pre-flow proposal, post-flow proposal, and unnormalized
one-step target density. \\
Flow maps and Jacobians &
$\pi_{t,\lambda}$, $v_{t,\lambda}$, $\Phi_t$, $J_t$, $D\Phi_t$ &
Homotopy density, pseudo-time velocity field, flow map, and forward Jacobian of
the flow. \\
Transport couplings &
$\Pi$, $\Pi_0^\star$, $\Pi_\varepsilon^\star$, $\Pi_{\varepsilon,K}$ &
Unregularized, entropy-regularized, and finite-Sinkhorn transport couplings. \\
Sinkhorn scalings &
$C$, $K_\varepsilon$, $a$, $b$, $\varepsilon$, $K$ &
Cost matrix, Gibbs kernel, row/column scalings, regularization level, and
finite iteration count. \\
Learned maps &
$T_\psi$, $Y_\tau$, $R_{\psi,\tau}$, $\mathcal D$ &
Learned student map, teacher output, student-teacher residual, and training
distribution. \\
HMC scalar and gradient &
$\ell_\star(u)$, $U_\star(u)$, $g(u)$, $H_\star(u,p)$ &
Named log target, potential, gradient used by the integrator, and Hamiltonian.
The value and gradient must refer to the same scalar. \\
\bottomrule
\end{longtable}
\end{center}

\begin{center}
\scriptsize
\begin{longtable}{@{}p{0.20\textwidth} p{0.32\textwidth} p{0.36\textwidth}@{}}
\caption{DPF-block claim-status vocabulary.}
\label{tab:bf-dpf-claim-status-index}\\
\toprule
Status phrase & Use only when & Disallowed shortcut \\
\midrule
\endfirsthead
\toprule
Status phrase & Use only when & Disallowed shortcut \\
\midrule
\endhead
Exact model identity &
The statement follows from the specified state-space model, Bayes' rule,
change of variables, or a named exact special case. &
Calling a closure, relaxation, solver, or learned map exact because it is
written in equations. \\
Unbiased particle estimator &
The expectation is taken over the stated particle randomness and the target
object is the likelihood or normalizing constant. &
Inferring unbiased log likelihood, score, pathwise gradient, or HMC validity. \\
Consistent approximation &
A convergence regime and target object are stated. &
Treating increasing $N$ or decreasing tolerance as a finite-sample proof. \\
Approximate closure &
A Gaussian, local-linear, finite-particle, or numerical closure replaces the
reference filtering object. &
Delaying the approximation label until a later chapter. \\
Relaxed target &
A smooth resampling or transport object replaces the classical categorical
operation and the scalar target is the relaxed object. &
Implying the original posterior is preserved without a correction. \\
Learned surrogate &
A trained map approximates a named teacher under a declared training
distribution. &
Using runtime, smoothness, or training loss as posterior validity evidence. \\
Engineering hypothesis &
The claim is a bounded recommendation or diagnostic rule, with required tests
stated. &
Presenting implementation preference as mathematical validation. \\
Unsupported claim &
The statement lacks a local derivation, reviewed source support, or documented
test evidence. &
Leaving the statement in the main argument without weakening or removing it. \\
\bottomrule
\end{longtable}
\end{center}

\section{Nonlinear state-space model and filtering recursion}
\label{sec:bf-pf-ssm-recursion}

Let $x_t \in \mathbb{R}^{n_x}$ denote the latent state and
$y_t \in \mathbb{R}^{n_y}$ the observation at time $t$.  Fix a parameter vector
$\theta \in \Theta \subset \mathbb{R}^{d_\theta}$.  The nonlinear state-space
model is specified by a transition density and an observation density,
\begin{align}
\label{eq:bf-pf-transition-density}
  x_t \mid x_{t-1} &\sim p_\theta(x_t \mid x_{t-1}), \\
\label{eq:bf-pf-observation-density}
  y_t \mid x_t &\sim p_\theta(y_t \mid x_t).
\end{align}
Write $y_{1:t}=(y_1,\ldots,y_t)$ for the observation history.  The predictive
law and the filtering law satisfy the standard Bayesian recursion
\begin{align}
\label{eq:bf-pf-predictive-law}
  p_\theta(x_t \mid y_{1:t-1})
  &= \int p_\theta(x_t \mid x_{t-1})
     p_\theta(x_{t-1} \mid y_{1:t-1})\,dx_{t-1}, \\
\label{eq:bf-pf-filtering-law}
  p_\theta(x_t \mid y_{1:t})
  &= \frac{p_\theta(y_t \mid x_t)
           p_\theta(x_t \mid y_{1:t-1})}
          {p_\theta(y_t \mid y_{1:t-1})},
\end{align}
where the one-step predictive density of the observation is
\begin{equation}
\label{eq:bf-pf-predictive-observation}
  p_\theta(y_t \mid y_{1:t-1})
  = \int p_\theta(y_t \mid x_t)
      p_\theta(x_t \mid y_{1:t-1})\,dx_t.
\end{equation}

The prediction equation is the Chapman--Kolmogorov identity applied to the
conditional law of $x_{t-1}$ given the observations already seen.  For any
measurable set $B$,
\[
  \mathbb{P}_\theta(x_t\in B\mid y_{1:t-1})
  =
  \int \mathbb{P}_\theta(x_t\in B\mid x_{t-1})
       p_\theta(x_{t-1}\mid y_{1:t-1})\,dx_{t-1}.
\]
When the transition kernel has density $p_\theta(x_t\mid x_{t-1})$, this set
identity gives \eqref{eq:bf-pf-predictive-law}.  The update equation is Bayes'
rule with the predictive law as prior and the observation density as
likelihood:
\[
  p_\theta(x_t\mid y_{1:t})
  =
  \frac{p_\theta(y_t\mid x_t,y_{1:t-1})
        p_\theta(x_t\mid y_{1:t-1})}
       {p_\theta(y_t\mid y_{1:t-1})}.
\]
The conditional-independence convention of the state-space model reduces
$p_\theta(y_t\mid x_t,y_{1:t-1})$ to $p_\theta(y_t\mid x_t)$, giving
\eqref{eq:bf-pf-filtering-law}.  Finally, repeated use of the chain rule gives
the marginal likelihood factorization
\[
  p_\theta(y_{1:T})
  =
  p_\theta(y_1)\prod_{t=2}^T p_\theta(y_t\mid y_{1:t-1}),
\]
with the same notation covering $t=1$ by interpreting $y_{1:0}$ as the empty
history.
The marginal likelihood then factorizes as
\begin{equation}
\label{eq:bf-pf-marginal-factorization}
  p_\theta(y_{1:T})
  = \prod_{t=1}^T p_\theta(y_t \mid y_{1:t-1}).
\end{equation}

Equations \eqref{eq:bf-pf-predictive-law}--\eqref{eq:bf-pf-marginal-factorization}
are exact.  The difficulty is computational rather than conceptual: outside
special classes such as the linear-Gaussian case, the integrals are rarely
available in closed form.  Particle filters replace these exact distributions by
random empirical approximations whose quality improves with particle count under
appropriate assumptions.

\section{Empirical filtering measures}
\label{sec:bf-pf-empirical-measures}

A particle filter represents the filtering law by a weighted empirical measure.
At time $t$, let
$\{(x_t^{(i)},w_t^{(i)})\}_{i=1}^N$ be particles and normalized weights, with
$w_t^{(i)}\ge 0$ and $\sum_{i=1}^N w_t^{(i)}=1$.  The empirical filtering
measure is
\begin{equation}
\label{eq:bf-pf-empirical-measure}
  \hat\pi_t^N(dx)
  = \sum_{i=1}^N w_t^{(i)}\,\delta_{x_t^{(i)}}(dx),
\end{equation}
where $\delta_z$ is the Dirac point mass at $z$.  The approximation of a test
function $\varphi$ under the filtering law is then
\begin{equation}
\label{eq:bf-pf-test-function-approx}
  \int \varphi(x)\,\hat\pi_t^N(dx)
  = \sum_{i=1}^N w_t^{(i)}\varphi(x_t^{(i)}).
\end{equation}
Thus the filter approximates the entire posterior law by a finite atomic
measure, not only its first two moments.

This is the main structural distinction between particle methods and the
Gaussian-approximation filters of the previous chapters.  The particle filter is
not built by closing the recursion in a finite-dimensional moment family.
Instead, it propagates a random discrete approximation to the law itself.

\section{Sequential importance sampling}
\label{sec:bf-pf-sis}

To derive the basic particle recursion, introduce a proposal density
$q_\theta(x_t \mid x_{0:t-1}, y_{1:t})$ from which particles can be sampled.
Assume that the proposal is positive whenever the trajectory target below is
positive.  Without this support condition the importance ratio is not a valid
correction; particles may miss regions that carry target probability.
Suppose a full trajectory proposal factors as
\begin{equation}
\label{eq:bf-pf-trajectory-proposal}
  q_\theta(x_{0:t} \mid y_{1:t})
  = q_\theta(x_0 \mid y_1)
    \prod_{s=1}^t q_\theta(x_s \mid x_{0:s-1}, y_{1:s}).
\end{equation}
The target filtering distribution over trajectories is proportional to
\begin{equation}
\label{eq:bf-pf-trajectory-target}
  p_\theta(x_{0:t}, y_{1:t})
  = p_\theta(x_0)
    \prod_{s=1}^t p_\theta(x_s \mid x_{s-1})
    p_\theta(y_s \mid x_s).
\end{equation}
The unnormalized importance weight is therefore
\begin{equation}
\label{eq:bf-pf-importance-weight}
  \tilde w_t(x_{0:t})
  = \frac{p_\theta(x_{0:t},y_{1:t})}
         {q_\theta(x_{0:t} \mid y_{1:t})}.
\end{equation}
This ratio is a trajectory-level object.  It is not yet a resampling rule and
not yet a differentiable likelihood map.  It is the Radon--Nikodym correction
that converts expectations under the proposal path law into expectations under
the unnormalized trajectory target, provided the support and integrability
conditions above hold.
Using the factorization in
\eqref{eq:bf-pf-trajectory-proposal}--\eqref{eq:bf-pf-trajectory-target}, one
obtains the recursive update
\begin{equation}
\label{eq:bf-pf-sis-recursion}
  \tilde w_t^{(i)}
  = \tilde w_{t-1}^{(i)}
    \frac{p_\theta(y_t \mid x_t^{(i)})
          p_\theta(x_t^{(i)} \mid x_{t-1}^{(i)})}
         {q_\theta(x_t^{(i)} \mid x_{0:t-1}^{(i)}, y_{1:t})}.
\end{equation}
Normalizing gives
\begin{equation}
\label{eq:bf-pf-normalized-weights}
  w_t^{(i)}
  = \frac{\tilde w_t^{(i)}}{\sum_{j=1}^N \tilde w_t^{(j)}}.
\end{equation}

Equation \eqref{eq:bf-pf-sis-recursion} is the generic sequential importance
sampling (SIS) update.  It already shows the two main mathematical stresses of
particle filtering:
\begin{enumerate}[label=(\roman*)]
  \item the proposal must be chosen carefully to avoid explosive weight
    variability;
  \item even when the estimator is well defined, the normalized weights tend to
    concentrate over time, which leads to degeneracy.
\end{enumerate}

\section{Sequential importance resampling and the bootstrap filter}
\label{sec:bf-pf-sir-bootstrap}

Sequential importance sampling alone is typically unstable over long horizons,
because the weights become increasingly uneven.  Sequential importance
resampling (SIR) introduces a resampling step to refresh the cloud when the
weight distribution becomes too concentrated.

The most classical special case is the bootstrap particle filter, which uses the
transition prior itself as the proposal:
\begin{equation}
\label{eq:bf-pf-bootstrap-proposal}
  q_\theta(x_t \mid x_{0:t-1}, y_{1:t})
  = p_\theta(x_t \mid x_{t-1}).
\end{equation}
Substituting \eqref{eq:bf-pf-bootstrap-proposal} into
\eqref{eq:bf-pf-sis-recursion} gives the simplified incremental weight
\begin{equation}
\label{eq:bf-pf-bootstrap-weight}
  \tilde w_t^{(i)}
  \propto w_{t-1}^{(i)} p_\theta(y_t \mid x_t^{(i)}).
\end{equation}
The bootstrap recursion at time $t$ can therefore be written as:
\begin{enumerate}[label=\textbf{Step \arabic*:}]
  \item \textbf{Select ancestors:}
    if resampling is performed before propagation, draw $a_{t-1}^{(i)}$ from
    the current normalized weights; otherwise keep $a_{t-1}^{(i)}=i$.
  \item \textbf{Propagate:}
    sample
    $x_t^{(i)} \sim p_\theta(x_t \mid x_{t-1}^{(a_{t-1}^{(i)})})$.
  \item \textbf{Weight:}
    compute
    $\tilde w_t^{(i)} = p_\theta(y_t \mid x_t^{(i)})$
    and normalize the weights.
  \item \textbf{Resample:}
    after recording the likelihood factor, either resample using the normalized
    weights and reset weights to $1/N$, or carry the weighted cloud forward
    until the next resampling decision.
\end{enumerate}
This yields the standard bootstrap/SIR particle filter of
\citet{gordon1993novel}, elaborated in the SMC text of
\citet{doucet2001sequential}.  The particle-MCMC literature
\citep{andrieu2009pseudo,andrieu2010particle} later uses the likelihood
estimator generated by this construction, but that pseudo-marginal role is a
statement about an unbiased value-side estimator inside a Metropolis--Hastings
construction.  It is not a statement that the realized resampling path is
smooth in $\theta$.

\section{The bootstrap particle-filter likelihood estimator}
\label{sec:bf-pf-likelihood-estimator}

The bootstrap particle filter supplies a natural estimator of the factors in the
marginal-likelihood decomposition
\eqref{eq:bf-pf-marginal-factorization}.  Define the average incremental weight
at time $t$ by
\begin{equation}
\label{eq:bf-pf-average-weight}
  Z_t^N = \frac{1}{N}\sum_{i=1}^N \tilde w_t^{(i)}.
\end{equation}
The corresponding likelihood estimator is
\begin{equation}
\label{eq:bf-pf-bootstrap-likelihood-estimator}
  \hat p_\theta^N(y_{1:T}) = \prod_{t=1}^T Z_t^N.
\end{equation}
The estimator-status result BayesFilter needs at this stage is the following.

\begin{proposition}[Bootstrap particle-filter likelihood-estimator status]
\label{prop:bf-pf-bootstrap-likelihood-status}
Fix the observations $y_{1:T}$ and the parameter $\theta$.  Assume the model
densities are dominated and integrable, the bootstrap propagation step samples
from $p_\theta(x_t\mid x_{t-1})$, and the Feynman--Kac quantities used below
are finite: for every bounded measurable $\varphi$ used in the induction,
the operator $Q_t\varphi$ is well defined, the incremental weights have finite
conditional expectation, and the displayed conditional expectations are finite.
Assume also that every resampling step used by the algorithm is conditionally
unbiased for the current normalized empirical measure.  Here conditional
unbiasedness means that, for every bounded measurable test function $\psi$ in
the same finite-expectation class, the expected post-resampling empirical
average equals the pre-resampling weighted empirical average, conditional on
the current weighted cloud.  Biased or deterministic approximation schemes that
do not satisfy this identity are not covered by the proposition.  Then the
estimator
\eqref{eq:bf-pf-bootstrap-likelihood-estimator} is an unbiased estimator of the
marginal likelihood:
\[
  \mathbb{E}\bigl[\hat p_\theta^N(y_{1:T})\bigr] = p_\theta(y_{1:T}).
\]
This claim is about the likelihood itself.  It does not claim that
$\log \hat p_\theta^N(y_{1:T})$ is an unbiased estimator of
$\log p_\theta(y_{1:T})$, and it does not claim that differentiating a realized
particle system gives an unbiased score.
\end{proposition}

\begin{proof}
For compactness, define the bootstrap Feynman--Kac operator
\[
  (Q_t\varphi)(x_{t-1})
  =
  \int \varphi(x_t)
       p_\theta(y_t\mid x_t)
       p_\theta(x_t\mid x_{t-1})\,dx_t
\]
for bounded measurable test functions $\varphi$.  Also define the unnormalized
filtering functional
\[
  \gamma_t(\varphi)
  =
  p_\theta(y_{1:t})
  \int \varphi(x_t)p_\theta(x_t\mid y_{1:t})\,dx_t .
\]
The prediction/update recursion implies
$\gamma_t(\varphi)=\gamma_{t-1}(Q_t\varphi)$ and
$\gamma_t(1)=p_\theta(y_{1:t})$.

The filtration convention is as follows.  At the end of time $t-1$ the
algorithm has a weighted empirical measure $\hat\pi_{t-1}^N$ and a
normalizing-constant estimate $\prod_{s=1}^{t-1}Z_s^N$.  Any resampling
randomness used to choose the particles propagated at time $t$ is then drawn,
and the transition propagation at time $t$ is conditionally independent across
particles given that post-resampling cloud.  All test functions in the induction
are bounded measurable, or otherwise integrable enough that the conditional
expectations below are finite.

Let $\mathcal F_{t-1}^N$ denote the sigma-field generated by the particle
system just before propagation at time $t$, after any resampling decision at
time $t-1$.  Let $\bar\pi_{t-1}^N$ be the empirical measure of the particles
from which the transition propagation is drawn.  If resampling is used, the
assumption of conditional unbiasedness means
\[
  \mathbb E[\bar\pi_{t-1}^N(\psi)\mid \hat\pi_{t-1}^N]
  =
  \hat\pi_{t-1}^N(\psi)
\]
for bounded $\psi$; if resampling is not used, then
$\bar\pi_{t-1}^N=\hat\pi_{t-1}^N$.

After bootstrap propagation and weighting,
\[
  Z_t^N \hat\pi_t^N(\varphi)
  =
  \frac{1}{N}\sum_{i=1}^N
  p_\theta(y_t\mid x_t^{(i)})\varphi(x_t^{(i)}).
\]
Conditional on $\mathcal F_{t-1}^N$, the propagated particles are drawn from
the transition kernel, so
\[
  \mathbb E\!\left[
    Z_t^N\hat\pi_t^N(\varphi)\mid \mathcal F_{t-1}^N
  \right]
  =
  \bar\pi_{t-1}^N(Q_t\varphi).
\]
Taking the conditional expectation over any resampling randomness replaces
$\bar\pi_{t-1}^N$ by $\hat\pi_{t-1}^N$ in expectation.  Multiplying by the
previous normalizing-constant estimate, which is measurable with respect to the
previous particle history, gives
\[
  \mathbb E\!\left[
    \left(\prod_{s=1}^t Z_s^N\right)\hat\pi_t^N(\varphi)
  \right]
  =
  \mathbb E\!\left[
    \left(\prod_{s=1}^{t-1} Z_s^N\right)
    \hat\pi_{t-1}^N(Q_t\varphi)
  \right].
\]
Iterating this identity and using
$\gamma_t(\varphi)=\gamma_{t-1}(Q_t\varphi)$ gives
\[
  \mathbb E\!\left[
    \left(\prod_{s=1}^t Z_s^N\right)\hat\pi_t^N(\varphi)
  \right]
  =
  \gamma_t(\varphi).
\]
Taking $\varphi\equiv 1$ yields
\[
  \mathbb{E}\left[\prod_{s=1}^t Z_s^N\right]
  =
  \gamma_t(1)
  =
  p_\theta(y_{1:t}).
\]
Setting $t=T$ proves the displayed likelihood-unbiasedness identity.  This is
the local BayesFilter version of the standard SMC normalizing-constant
unbiasedness argument; see \citet{doucet2001sequential} and
\citet{andrieu2010particle} for the full particle-system treatment.
\end{proof}

This proposition is mathematically central for the later DPF program.  It tells
us that the classical bootstrap particle filter supplies an unbiased estimator
of the marginal likelihood, even though it does not supply a pathwise differentiable
likelihood map.  Later differentiable constructions must therefore be compared
against this baseline carefully: a differentiable surrogate may be smoother to
optimize or differentiate, but it need not preserve this exact estimator
status.

\section{Differentiability boundary}
\label{sec:bf-pf-differentiability-boundary}

The unbiasedness result above is a value-side Monte Carlo statement.  It should
not be read as a pathwise differentiability statement.  For a fixed random seed,
the propagation step may be reparameterizable in some models, but the
resampling step is a categorical map from weights to ancestor indices.  Small
changes in $\theta$ can leave the selected ancestor unchanged over an interval
and then switch it discontinuously when a cumulative weight boundary is crossed.
The realized resampled cloud is therefore not a smooth function of $\theta$.

This distinction creates three separate objects that later chapters must not
blur:
\begin{enumerate}[label=(\roman*)]
  \item the true likelihood $p_\theta(y_{1:T})$;
  \item an unbiased randomized estimator
    $\hat p_\theta^N(y_{1:T})$ of that likelihood;
  \item a pathwise differentiable scalar produced by a relaxed, transported, or
    learned particle computation.
\end{enumerate}
The first object is the statistical target, the second can support
pseudo-marginal value-side constructions under their own assumptions, and the
third may be useful for optimization or HMC only after its target status is
stated.  A gradient through the third object is not automatically a gradient of
the first or of the expectation of the second.

\section{Degeneracy and effective sample size}
\label{sec:bf-pf-degeneracy}

The main numerical weakness of the particle filter is weight degeneracy.  As the
observation sequence grows, the normalized weights frequently concentrate on a
small subset of particles.  A standard summary is the effective sample size
(ESS),
\begin{equation}
\label{eq:bf-pf-ess}
  \mathrm{ESS}_t = \frac{1}{\sum_{i=1}^N (w_t^{(i)})^2}.
\end{equation}
The extreme cases are informative:
\begin{itemize}
  \item if all weights are equal, then $\mathrm{ESS}_t=N$;
  \item if one particle carries almost all the mass, then
    $\mathrm{ESS}_t \approx 1$.
\end{itemize}
When ESS collapses, the atomic approximation
\eqref{eq:bf-pf-empirical-measure} still exists, but its practical variance may
become extremely large.  Resampling alleviates this by discarding low-weight
particles and duplicating high-weight particles, but resampling does not remove
the deeper dimensionality problem.  The classical literature shows that in many
high-dimensional settings the number of particles needed to stabilize the weight
distribution can grow very rapidly with the effective dimension of the
observation problem \citep{bengtsson2008curse,snyder2008obstacles}.

This is the exact point where the later chapters enter.  Particle-flow methods,
PF-PF corrections, differentiable resampling, OT relaxations, and learned OT
operators are all responses to some combination of:
\begin{enumerate}[label=(\roman*)]
  \item degeneracy under standard resampling;
  \item lack of differentiability of the resampling step;
  \item computational burden of high-quality resampling or transport maps;
  \item mismatch between a useful learning objective and a valid HMC target.
\end{enumerate}
A mathematically clean DPF monograph must therefore begin here: by making the
classical particle-filter object precise before modifying it.

\section{Baseline audit table}
\label{sec:bf-pf-audit-table}

\begin{center}
\scriptsize
\begin{longtable}{@{}p{0.20\textwidth} p{0.23\textwidth} p{0.24\textwidth} p{0.23\textwidth}@{}}
\caption{Classical particle-filter baseline audit.}
\label{tab:bf-pf-baseline-audit}\\
\toprule
Object or claim & Status & Assumptions & Implementation diagnostic \\
\midrule
\endfirsthead
\toprule
Object or claim & Status & Assumptions & Implementation diagnostic \\
\midrule
\endhead
Filtering recursion and likelihood factorization &
Exact model identity &
Dominated transition and observation kernels; fixed $\theta$; finite
normalizers. &
Compare predictive and filtering factors against a Kalman reference in a
linear-Gaussian model. \\
Trajectory importance ratio &
Exact algebraic identity for the specified target/proposal pair &
Proposal support covers target support; finite weights. &
Check target and proposal log densities separately on a small trajectory. \\
Weighted empirical measure &
Finite-$N$ Monte Carlo approximation &
Particle system generated by the declared proposal and weighting rule. &
Compare bounded test-function expectations across increasing $N$. \\
Bootstrap likelihood estimator &
Unbiased estimator of likelihood, not log likelihood or score &
Bootstrap proposal; conditionally unbiased resampling; integrable likelihood
factors. &
In a controlled model, average repeated estimates and compare to the analytic
likelihood. \\
ESS and ancestor diagnostics &
Engineering diagnostic, not posterior-validity proof &
Weights normalized and finite; diagnostic interpreted with model dimension and
observation informativeness. &
Track ESS, log-weight variance, and unique ancestor count by time and
parameter region. \\
Pathwise gradient through resampling &
Not supplied by the classical bootstrap filter &
Categorical ancestor map is used without relaxation. &
Finite-difference a fixed-seed resampled path and identify discontinuities or
zero-almost-everywhere regions. \\
\bottomrule
\end{longtable}
\end{center}

\section{What this chapter does and does not claim}
\label{sec:bf-pf-chapter-boundary}

This chapter establishes the exact nonlinear filtering recursion and the
classical bootstrap particle-filter baseline.  It does \emph{not} yet claim any
of the following:
\begin{itemize}
  \item a pathwise differentiable likelihood map;
  \item a particle-flow approximation;
  \item a proposal-corrected flow filter;
  \item a differentiable resampling rule;
  \item an HMC-ready particle backend.
\end{itemize}
Those belong to the later chapters.  The role of the present chapter is more
basic and more important: it fixes the reference object against which every
later DPF construction must be interpreted.
